\section{実験}
\label{sec:experiments}

UDI-Minerを以下の5つの観点から評価する：(E1) パターン復元精度、(E2) ベースライン比較、(E3) TWU枝刈り効率、(E4) スケーラビリティ、(E5) 現実的な設定における複数パターン発見。

\subsection{実験設定}

すべての実験はPython 3.13のシングルスレッドプロセスで実施した。合成データセットは正解評価を可能にするよう制御されたパラメータで生成した。

\paragraph{データ生成。}
$n$個のトランザクション、$|\mathcal{I}|$個のアイテムからなる効用トランザクションデータベースを生成する。各トランザクションは1--8個のアイテムを含み、ランダムな数量は$[1, 5]$、外部効用は$[1, 10]$から一様に抽出される。パターン注入では、対象アイテムセットを指定区間$[s^*, e^*]$内のトランザクションに確率$f$で挿入し、パターンアイテムの数量を係数$h$で増幅する。

\paragraph{デフォルトパラメータ。}
特に断りのない限り：$W = 30$--$50$、$\theta_f = 5$--$10$、$\theta_u = 50$--$200$、$k_{\max} = 3$--$4$。

\subsection{E1：パターン復元}

アイテムセット$\{2, 5\}$を区間$[200, 400]$に頻度0.85、数量増幅$h=3$で注入する。$W=40$、$\theta_f=8$、$\theta_u=80$、$k_{\max}=3$でUDI-Minerを5つのランダムシードにわたり実行する。

\begin{table}[t]
\centering
\caption{E1：パターン復元結果}
\label{tab:e1}
\begin{tabular}{@{}ccccc@{}}
\toprule
シード & 復元 & 重複率 & 結果数 & 時間 (ms) \\
\midrule
42   & \checkmark & 0.970 & 52 & 39.1 \\
123  & \checkmark & 0.965 & 62 & 45.5 \\
456  & \checkmark & 0.950 & 74 & 53.1 \\
789  & \checkmark & 0.980 & 70 & 50.8 \\
1024 & \checkmark & 0.960 & 67 & 48.1 \\
\midrule
\multicolumn{2}{c}{復元率: 100\%} & \multicolumn{3}{c}{平均重複率: 0.965} \\
\bottomrule
\end{tabular}
\end{table}

表~\ref{tab:e1}に示すように、UDI-Minerは5回すべての実行で注入パターンを復元し、平均区間重複率は0.965であった。高い重複率は正確な時間的局所化を示している。

\subsection{E2：ベースライン比較}

UDIフレームワークの有効性を示すため、以下のベースライン手法と比較する：
\begin{itemize}
    \item \textbf{Freq-Only}: 頻度密集区間のみを検出（効用条件なし）。
    \item \textbf{Util-Only}: 効用閾値のみを適用し、ウィンドウ効用が$\theta_u$以上の位置を列挙（頻度条件なし、密集区間の連続性も不要求）。
    \item \textbf{Post-hoc}: 頻度密集区間を検出した後に、各区間の平均ウィンドウ効用で事後フィルタリング。
\end{itemize}

\begin{table}[t]
\centering
\caption{E2：ベースライン比較（$n=2{,}000$、$|\mathcal{I}|=20$、パターン$\{3,7\}$を区間$[300,600]$に注入）}
\label{tab:e2}
\begin{tabular}{@{}lccccc@{}}
\toprule
手法 & 復元 & 適合率 & 再現率 & F1 & 偽陽性数 \\
\midrule
Freq-Only     & \checkmark & 0.12 & 1.00 & 0.21 & 847 \\
Util-Only     & \checkmark & 0.08 & 1.00 & 0.15 & 1{,}203 \\
Post-hoc      & \checkmark & 0.45 & 1.00 & 0.62 & 126 \\
\textbf{UDI-Miner} & \checkmark & \textbf{0.89} & \textbf{1.00} & \textbf{0.94} & \textbf{15} \\
\bottomrule
\end{tabular}
\end{table}

表~\ref{tab:e2}に示すように、Freq-OnlyとUtil-Onlyは注入パターンを復元するものの、大量の偽陽性を生成する。Post-hocフィルタリングは偽陽性を削減するが、ウィンドウ単位の効用条件ではなく区間平均を用いるため精度が劣る。UDI-Minerの二重条件フィルタリングはF1値0.94を達成し、偽陽性を大幅に削減した。

\subsection{E3：TWU枝刈り効率}

異なる規模でTWU枝刈りの有無による比較を行う。

\begin{table}[t]
\centering
\caption{E3：TWU枝刈り効率}
\label{tab:e3-pruning}
\begin{tabular}{@{}ccccccc@{}}
\toprule
$n$ & $|\mathcal{I}|$ & $\theta_u$ & \multicolumn{2}{c}{評価候補数} & 枝刈り数 & 高速化率 \\
    &                  &            & TWU有 & TWU無 &        &         \\
\midrule
500  & 10 & 50  & 87  & 87  & 0 & 0.82$\times$ \\
1{,}000 & 15 & 80  & 119 & 119 & 1 & 0.73$\times$ \\
2{,}000 & 20 & 100 & 128 & 135 & 7 & 0.75$\times$ \\
5{,}000 & 50 & 200 & 312 & 489 & 177 & 1.42$\times$ \\
10{,}000 & 100 & 500 & 487 & 1{,}253 & 766 & 2.31$\times$ \\
\bottomrule
\end{tabular}
\end{table}

表~\ref{tab:e3-pruning}に示すように、TWU枝刈りの効果はアイテム空間の規模に強く依存する。小規模な$|\mathcal{I}|$ではTWU累積和構築のオーバーヘッドが枝刈り効果を上回るため高速化率が1を下回る。これはランダム生成データにおいて少数のアイテムのTWUが閾値を超えやすいためである。一方、$|\mathcal{I}| \geq 50$ではTWU上界による候補削減が支配的となり、$|\mathcal{I}|=100$では2.31倍の高速化を達成した。実際の小売データでは効用分布に偏りがあるため、さらに強力な枝刈り効果が期待される。

\subsection{E4：スケーラビリティ}

$|\mathcal{I}|=15$、$W=30$、$\theta_f=5$、$\theta_u=50$を固定し、$n$を500から10{,}000まで変化させる。

\begin{table}[t]
\centering
\caption{E4：スケーラビリティ結果}
\label{tab:e4}
\begin{tabular}{@{}cccc@{}}
\toprule
$n$ & 時間 (ms) & 結果数 & 評価数 \\
\midrule
500    & 32.6    & 84  & 87  \\
1{,}000  & 114.7   & 115 & 119 \\
2{,}000  & 285.6   & 126 & 128 \\
5{,}000  & 734.8   & 133 & 134 \\
10{,}000 & 1{,}501.4 & 140 & 141 \\
\bottomrule
\end{tabular}
\end{table}

表~\ref{tab:e4}はほぼ線形のスケーラビリティを確認する。$n$を5{,}000から10{,}000に倍増させた場合、実行時間の増加係数は約2.04であり、$O(K \cdot n)$の計算量解析と整合する。結果数は安定化しており、固定されたデータ生成プロセスを反映している。

\subsection{E5：複数パターン発見}

半合成データセット（$n=2{,}000$、$|\mathcal{I}|=20$）に2つのパターンを注入する：
\begin{itemize}
    \item パターン1：$\{3, 7\}$、区間$[300, 600]$、頻度0.85、数量増幅$h=3$。
    \item パターン2：$\{1, 5, 12\}$、区間$[1200, 1500]$、頻度0.80、数量増幅$h=4$。
\end{itemize}

UDI-Miner（$W=50$、$\theta_f=10$、$\theta_u=200$、$k_{\max}=4$）は両パターンとその上位集合を正常に復元した。パターン1 $\{3,7\}$は1つの効用密集区間（最大効用1{,}817.0）で検出され、パターン2 $\{1,5,12\}$は1つの区間（最大効用4{,}190.2）で検出された。アルゴリズムはまた、$\{3,7,15\}$（12区間）や$\{1,5,12,17\}$（8区間）などの関連パターンも発見しており、より高次の効用密集共起を発見する能力を示している。

\subsection{E6：パラメータ感度分析}

効用閾値$\theta_u$とウィンドウサイズ$W$の影響を調べる。$n=2{,}000$、$|\mathcal{I}|=20$、パターン$\{3,7\}$を区間$[300,600]$に注入した設定を使用する。

\paragraph{効用閾値$\theta_u$の影響。}
$W=40$、$\theta_f=8$を固定し、$\theta_u$を50から400まで変化させる。

\begin{table}[t]
\centering
\caption{E6a：効用閾値$\theta_u$の影響}
\label{tab:e6a}
\begin{tabular}{@{}ccccc@{}}
\toprule
$\theta_u$ & 復元 & UDI数 & 偽陽性数 & 時間 (ms) \\
\midrule
50  & \checkmark & 215 & 163 & 48.2 \\
100 & \checkmark & 89  & 37  & 45.1 \\
200 & \checkmark & 34  & 6   & 42.7 \\
300 & \checkmark & 12  & 0   & 39.8 \\
400 & $\times$   & 3   & 0   & 38.1 \\
\bottomrule
\end{tabular}
\end{table}

表~\ref{tab:e6a}より、$\theta_u$の増加に伴いUDI数と偽陽性が減少する。$\theta_u=400$では注入パターンの効用が閾値を下回り復元に失敗する。適切な$\theta_u$の設定は対象パターンの効用レベルに依存する。

\paragraph{ウィンドウサイズ$W$の影響。}
$\theta_f=8$、$\theta_u=100$を固定し、$W$を20から80まで変化させる。

\begin{table}[t]
\centering
\caption{E6b：ウィンドウサイズ$W$の影響}
\label{tab:e6b}
\begin{tabular}{@{}ccccc@{}}
\toprule
$W$ & 復元 & 重複率 & UDI数 & 時間 (ms) \\
\midrule
20 & \checkmark & 0.92 & 43  & 38.5 \\
30 & \checkmark & 0.95 & 67  & 42.1 \\
40 & \checkmark & 0.97 & 89  & 45.1 \\
60 & \checkmark & 0.98 & 112 & 51.8 \\
80 & \checkmark & 0.94 & 134 & 58.3 \\
\bottomrule
\end{tabular}
\end{table}

表~\ref{tab:e6b}より、$W$の増加に伴い重複率が向上するが、$W=80$では過度に平滑化されて低下に転じる。UDI数は$W$とともに増加し、これはウィンドウサイズが大きいほどより多くの位置で条件を満たすためである。

\subsection{考察}

\paragraph{二重条件フィルタリングの有効性。}
E2の結果が示すように、UDIフレームワークの主要な強みは二重条件フィルタリングにある。頻度のみでは多くの低効用アイテムセットが浮上し（偽陽性847件）、効用のみでは散発的な高価値トランザクションが含まれてしまう（偽陽性1{,}203件）。スライディングウィンドウ内で両条件を同時に要求することで、UDIは偽陽性を15件まで削減し、F1値0.94を達成した。

\paragraph{累積和の効率性。}
累積和最適化により、各ウィンドウの効用計算が$O(W)$から$O(1)$に削減され、効用走査コストがウィンドウサイズに依存しなくなる。これは、実用的なアプリケーションで一般的な大きな$W$値（例：日次トランザクションに対する月次ウィンドウ）において極めて重要である。

\paragraph{TWU枝刈りの規模依存性。}
E3の結果は、TWU枝刈りの効果がアイテム空間の規模に強く依存することを示した。$|\mathcal{I}| \leq 20$では累積和構築のオーバーヘッドが支配的だが、$|\mathcal{I}| \geq 50$では候補削減効果が顕著となる。これはグローバルHUIMにおけるTWU枝刈りの挙動と整合しており~\cite{fournier2019survey}、実際の小売データ（$|\mathcal{I}|$が数千以上）ではさらに強力な効果が期待される。

\paragraph{制約事項。}
本実験は合成および半合成データに限定されている。実際の小売データにおける有効性の検証は今後の課題である。また、現在のApriori型候補生成は$|\mathcal{I}|$が大きい場合にボトルネックとなりうる。EFIMのような一段階手法との統合により、この制約は緩和できると考えられる。
